\SourceLeaf{82}%
\begingroup\postdisplaypenalty=10000
\[
1+d+E+367(A-1700)-2b
\]
\endgroup
oder
\[
d+E+367(A-1700)-2b
\]
Dann ziehe ich $364(A-1700)$ ab. Das ergibt
\[
d+E+3A-5099-2b
\]
oder
\[
d+E+3A-5100-2b
\]
Addiert man anschließend 5096, so erhält man
\[
d+E+3A-3-2b
\]
oder
\[
d+E+3A-4-2b
\]
Zieht man schließlich $3A-3c$ ab, was offensichtlich durch 7 teilbar ist, so erhält man
\[
d+E+3c-3-2b
\]
oder
\[
d+E+3c-4-2b
\]
Zieht man dies von $7c+7d$ ab, so erhält man
\[
3+2b+4c+6d-E
\]
oder
\[
4+2b+4c+6d-E
\]
was also durch 7 teilbar sein muss. Daraus ist klar, dass $E$ der Rest ist, den man erhält, wenn man
\[
3+2b+4c+6d
\]
oder
\[
4+2b+4c+6d
\]
durch 7 teilt. Folglich ist $E=e$.

Ostern fällt also auf den $22+d+e^{\mathrm{ten}}$ März oder, was dasselbe ist, auf den $d+e-9^{\mathrm{ten}}$ April. W. Z. B. W.

\Needspace{13\baselineskip}
\SourceLeaf{83}%
\begin{center}\textit{Ganz allgemeine Regeln zur Berechnung des Osterfestes\\
nach dem julianischen und nach dem gregorianischen Kalender.}\end{center}
\begin{center}
\begin{tabular}{l|r|l}
Die Division & durch & ergibt den Rest\\\hline
der Jahreszahl & 19 & $a$\\
der Jahreszahl & 4 & $b$\\
der Jahreszahl & 7 & $c$\\
der Zahl $19a+M$ & 30 & $d$\\
der Zahl $2b+4c+6d+N$ & 7 & $e$\\\hline
\end{tabular}\par
Dann fällt Ostern auf den $22+d+e^{\mathrm{ten}}$ März\par
oder den $d+e-9^{\mathrm{ten}}$ April
\end{center}

$M$ und $N$ sind Zahlen, deren Werte im julianischen Kalender für immer, im gregorianischen dagegen jeweils für mindestens 100 Jahre unverändert bleiben. Im julianischen Kalender gilt $M=15$, $N=6$. Im gregorianischen gilt von seiner Einführung bis 1699 $M=22$, $N=2$.

\begin{center}\small
\setlength{\tabcolsep}{5pt}\begin{tabular}{rrll|rrll}
von 1700 & bis 1799 & $M=23$, & $N=3$ & von 2100 & bis 2199 & $M=24$, & $N=6$\\
1800 & $\ldots$1899 & $M=23$, & $N=4$ & 2200 & $\ldots$2299 & $M=25$, & $N=0$\\
1900 & $\ldots$1999 & $M=24$, & $N=5$ & 2300 & $\ldots$2399 & $M=26$, & $N=1$\\
2000 & $\ldots$2099 & $M=24$, & $N=5$ & 2400 & $\ldots$2499 & $M=25$, & $N=1$
\end{tabular}
\end{center}

Im Allgemeinen findet man im gregorianischen Kalender die Werte von $M$ und $N$ für jedes gegebene Jahrhundert von $100k$ bis $100k+99$ nach der folgenden Regel.

Es ergebe
\[
k\text{ durch }\left\{\begin{smallmatrix}3\\4\end{smallmatrix}\right\}
\text{ geteilt die (ganzzahligen) Quotienten }
\left\{\begin{smallmatrix}p\\q\end{smallmatrix}\right\}
\]
wobei die Reste unberücksichtigt bleiben.
Dann ist
\[
\left\{\begin{smallmatrix}M\\N\end{smallmatrix}\right\}
\text{ der \textit{Rest}, den man erhält, wenn man }
\left\{\begin{smallmatrix}15+k-p-q\\4+k-q\end{smallmatrix}\right\}
\text{ durch }\left\{\begin{smallmatrix}30\\7\end{smallmatrix}\right\}\text{ teilt.}
\]
\textit{Beispiel.} Für die 100 Jahre von 4700 bis 4799 ist $k=47$, $p=15$, $q=11$. Daraus ergibt sich $15+k-p-q=36$ und $4+k-q=40$, also $M=6$, $N=5$. So erhält man zum Beispiel für das Jahr 4763
\SourceLeaf{84}%
\begin{center}\small
\begin{tabular}{r|r|l}
$a=13$ & $19a+M=253$ & $e=3$\\
$b=3$ & $d=13$ & Ostern am $13+3-9$, das heißt am 7. April\\
$c=3$ & $2b+4c+6d+N=101$ & \quad nach dem greg. Kalender\\
\multicolumn{3}{@{}c@{}}{}\\[-4pt]
\multicolumn{2}{@{}l|}{\shortstack[l]{Nach dem julianischen\\Kalender dagegen}} & \\
\multicolumn{2}{r|}{$19a+M=262$} & $e=2$\\
\multicolumn{2}{r|}{$d=22$} & Ostern am $22+2-9$, das heißt am\\
\multicolumn{2}{r|}{$2b+4c+6d+N=156$} & \quad 15. April
\end{tabular}
\end{center}

Von den obigen Regeln gibt es im \textit{gregorianischen Kalender} nur die folgenden zwei Ausnahmen.

I. Ergibt die Rechnung den 26. April als Ostertermin, so nimmt man stattdessen \textit{immer} den 19. April. (Z.\,B. 1609. 1989).

Man sieht leicht, dass dieser Fall nur dann vorkommen kann, wenn die Rechnung $d=29$ und $e=6$ ergibt. Den Wert 29 kann $d$ nur dann annehmen, wenn bei der Division von $11M+11$ durch 30 ein Rest herauskommt, der \textit{kleiner} als 19 ist. Dazu muss $M$ einen der folgenden 19 Werte haben.
\[
\begin{gathered}
0,\ 2,\ 3,\ 5,\ 6,\ 8,\ 10,\ 11,\ 13,\ 14,\ 16,\ 17,\ 19,\ 21,\ 22,\ 24,\ 25,\ 27,\ 29
\end{gathered}
\]
II. Ergibt die Rechnung $d=28$, $e=6$, und ergibt die Division von $11M+11$ durch 30 zudem einen Rest, der kleiner als 19 ist, so ergibt die Rechnung den 25. April, Ostern fällt jedoch auf den 18. April. --- Man überzeugt sich leicht, dass dieser Fall nur in denjenigen Jahrhunderten eintreten kann, in denen $M$ einen der folgenden acht Werte hat.
\[
2,\ 5,\ 10,\ 13,\ 16,\ 21,\ 24,\ 29
\]
\textit{Von diesen beiden Ausnahmen abgesehen, gelten die obigen Regeln ganz allgemein.}
\par\bigskip\centerline{\rule{.24\textwidth}{.3pt}}\medskip
[Handschriftliche Bemerkung:]

$p$ wird als Quotient bei der Division von $8k+13$ durch 25 bestimmt.
\begin{center}\small
\setlength{\tabcolsep}{4pt}
\begin{tabular}{lrl}
Von Ostern bis Michaelis sind es im Durchschnitt $\ .\ .$ & $173\frac{52}{95}$ & Tage\\
und (Ostern eingeschlossen, Michaelis ausgeschlossen) $\ .\ .\ .$ & $25\frac{21}{95}$ & Sonntage\\
& $148\frac{31}{95}$ & Wochentage\\
Von Michaelis bis Ostern dagegen $\ .\ .\ .\ .\ .\ .$ & $164\frac{773}{950}$ & Wochentage
\end{tabular}
\end{center}
\par\bigskip\centerline{\rule{.24\textwidth}{.3pt}}
